\section{总结与延伸}

Lecture 1 建立了 CS336 的整体方法：

\begin{enumerate}
    \item LLM 已经工业化，frontier model 昂贵且不透明，所以课堂要训练可迁移 mechanics 和 mindset。
    \item 小模型能教实现和调试，但不能保证所有系统瓶颈和能力表现都能外推。
    \item 现代语言模型技术栈来自概率建模、神经表示、Transformer、系统并行、scaling laws、数据工程和 alignment 的叠加。
    \item Tokenizer 是第一块真正可实现的积木，它把文本变成模型可训练的离散序列。
\end{enumerate}

\begin{importantbox}{最终 takeaway}
From scratch 不是怀旧，而是研究方法。只有把 tokenizer、architecture、training loop、optimizer、hardware、data 和 inference 都拆开，才能知道模型质量来自哪里，也知道该从哪里改进。
\end{importantbox}

\subsection{拓展阅读}

\begin{itemize}
    \item Shannon 1950: entropy and prediction of English.
    \item Sennrich et al. 2016: Byte Pair Encoding for neural machine translation.
    \item Vaswani et al. 2017: Transformer.
    \item Kaplan et al. 2020 / Chinchilla 2022: scaling laws and compute-optimal training.
    \item JAX Scaling Book: systems and resource accounting for large models.
\end{itemize}

\subsection{Executable source 索引：每个课程节点落在哪里}

Lecture 1 的源材料不是传统 PDF deck，而是一份由 section、text、figure 与 code 节点组成的 executable lecture。下面的索引保留源节点原名，并说明讲义如何把它们组织成教学主线；它既是 coverage ledger，也提醒读者哪些标题只是导航，哪些节点在后文会展开成完整机制。

\begin{longtable}{p{0.43\textwidth}p{0.47\textwidth}}
\toprule
\textbf{Executable-source titles} & \textbf{讲义中的处理} \\
\midrule
CS336: Language Models From Scratch (Spring 2026)\newline
Why did we make this course? & 课程定位、from-scratch 方法、API 抽象边界与 executable lecture。 \\
The industrialization of language models\newline
What can we learn in this class that transfers to frontier models?\newline
Intuitions?\newline
The bitter lesson\newline
accuracy = efficiency x resources & 工业化成本、尺度不可直接外推、mechanics/mindset/intuitions 与可扩展算法。 \\
Pre-neural (before 2010s)\newline
Neural ingredients (2010s)\newline
Early foundation models (late 2010s)\newline
Embracing scaling\newline
Open models & 从 Shannon、n-gram 到 Transformer、scaling 与开放模型生态的历史主线。 \\
Why you should take this course\newline
Why you should not take this course\newline
How you can follow along at home\newline
Assignments\newline
Compute & 课程工作量、适合人群、在家复现方式、五份作业、AI policy 与 compute 支持。 \\
Model architecture\newline
Assignment 1 (basics)\newline
Byte Pair Encoding (BPE)\newline
Training the tokenizer\newline
Using the tokenizer & Basics 单元和 tokenizer 章节：模型组件、BPE 训练、encode/decode 与第一份作业。 \\
Kernels\newline
Parallelism\newline
Assignment 2 (systems) & Systems 单元：单 GPU kernel、多 GPU 并行、optimizer-state sharding 与性能验收。 \\
Assignment 3 (scaling laws)\newline
Evaluation & Scaling 单元：小规模拟合、预注册 loss 预测，以及不能只看 training loss 的评测边界。 \\
Data curation\newline
Data processing\newline
Assignment 4 (data) & Data 单元：来源选择、HTML/PDF 清洗、过滤、去重、混合与固定预算 leaderboard。 \\
Assignment 5 (alignment)\newline
Observations & Alignment 单元：偏好优化、RL systems，以及模型能力、行为目标和资源约束之间的观察。 \\
\bottomrule
\end{longtable}

\begin{knowledgebox}{如何使用这张源节点索引}
若只想理解第一讲，可按正文顺序阅读；若要回到 executable source 调试或扩写，就用左列标题定位节点，再检查同一节点附近的 text、figure 与 code。导航标题本身不提供完整解释，讲义中的表格、公式、读图框和 teacher-voice 段落才是这些节点的教学展开。
\end{knowledgebox}

\end{document}
